\section{Постановка задачи}\label{sec:problem}

\subsection{Данные, кандидаты и доступная информация}
Пусть $X=(x_1,\ldots,x_n)$, $x_i\in\mathbb R^d$, --- фиксированный набор объектов, а $y=(y_1,\ldots,y_n)$ --- скрытые метки целевых классов. До раскрытия меток строится конечный пул $\mathcal P=\{P_1,\ldots,P_M\}$. Элемент пула представляет полное разбиение набора. Кандидат $P$ задаёт принадлежность к кластеру для каждого объекта. Номера кластеров номинальны. Пул может включать результаты разных алгоритмов, чисел кластеров и настроек. Одинаковые разбиения объединяются.

Множество индексов объектов разделяется на резерв $S$ и проверочную часть $T$, $S\cap T=\varnothing$. В одном повторении из $S$ раскрывается подмножество $L$, $|L|=m$. Правило выбора получает признаки всех объектов, пул и метки $y_L$, но не получает $y_T$ или метки нераскрытой части резерва. Предобработка признаков без меток и построение кандидатов на всём $X$ разрешены трансдуктивной постановкой. Известные связанные объекты объединяются в группы и при разделении перемещаются совместно.

\subsection{Внутренняя и внешняя оценки}
Основной внутренний индекс $I(P)$ --- средний silhouette~\citep{rousseeuw1987}. Для точки $i$ обозначим через $a_i$ среднюю дистанцию до остальных объектов своего кластера, а через $b_i$ --- минимальную среднюю дистанцию до другого кластера. Тогда
\begin{equation}
 s_i=\frac{b_i-a_i}{\max(a_i,b_i)},\qquad I(P)=\frac1n\sum_{i=1}^n s_i.
\end{equation}
Для одноэлементного кластера и случая $a_i=b_i=0$ используется соглашение $s_i=0$. Дополнительная геометрическая оценка --- индекс Davies--Bouldin~\citep{davies1979}. При выборе направление его оптимизации меняется на противоположное, поскольку меньшее значение предпочтительно.

Частичная внешняя оценка и проверочное качество определяются на разных объектах:
\begin{equation}\label{eq:two-ari}
 A_L(P)=\ARI(y_L,P_L),\qquad Q_T(P)=\ARI(y_T,P_T).
\end{equation}
Пусть $n_{ab}$ --- число размеченных объектов класса $a$, попавших в кластер $b$. $n_{a\cdot}$ и $n_{\cdot b}$ --- суммы по строкам и столбцам таблицы сопряжённости. Для $m\ge2$ положим $U=\sum_{a,b}\binom{n_{ab}}2$, $V=\sum_a\binom{n_{a\cdot}}2$, $W=\sum_b\binom{n_{\cdot b}}2$, $N=\binom m2$. ARI имеет вид~\citep{hubert1985}
\begin{equation}
 \ARI=\frac{U-VW/N}{(V+W)/2-VW/N}.
\end{equation}
В вырожденных случаях используются соглашения реализации scikit-learn, в том числе значение 1 для совпадающих тривиальных разбиений. Такие случаи отдельно диагностируются.

Если алгоритм оставляет шумовые объекты без назначения, для внешних индексов каждый такой объект рассматривается как отдельная одноэлементная группа. В реализованном протоколе внутренние показатели сначала рассчитываются на назначенных объектах. Для доли назначенных $q$ приняты преобразования $S\mapsto qS$ и $DB\mapsto DB/q$. Кандидаты с $q<0{,}7$ или менее чем двумя кластерами исключаются. Эти соглашения фиксируются до сравнения методов. Преобразование $qS$ при отрицательном silhouette увеличивает значение индекса при уменьшении $q$, поэтому его нельзя интерпретировать как универсальный штраф за шум. Влияние обработки шумовых объектов требует отдельной проверки чувствительности.

\subsection{Правило выбора с адаптивным допуском}
Внешний лидер определяется только по доступной разметке:
\begin{equation}\label{eq:leader}
 \Pari=\arg\max_{P\in\mathcal P}A_L(P).
\end{equation}
Равенства разрешаются фиксированным порядком канонических идентификаторов кандидатов. Для каждого $P$ вычисляется наблюдаемое отставание
\begin{equation}
 d_P=A_L(\Pari)-A_L(P).
\end{equation}
Из $L$ формируются $B$ бутстрепных выборок $L^{(1)},\ldots,L^{(B)}$: объекты выбираются с возвращением, а при наличии связанных объектов перевыбираются целые группы. Повторения объектов учитываются с их кратностями. Пары объектов независимо не перевыбираются. Кластеризации и внутренние индексы остаются фиксированными.

Для всех кандидатов используются одни и те же бутстрепные выборки. Лидер из~\eqref{eq:leader} фиксируется на исходном $L$ и не переопределяется при вычислении
\begin{equation}
 d_P^{(b)}=A_{L^{(b)}}(\Pari)-A_{L^{(b)}}(P),\qquad
 \widehat\sigma_P=\sd_{b=1}^B(d_P^{(b)}).
\end{equation}
При $B\ge2$ оценка разброса задаётся явно:
\begin{equation}
 \bar d_P=\frac1B\sum_{b=1}^B d_P^{(b)},\qquad
 \widehat\sigma_P=\sqrt{\frac1{B-1}\sum_{b=1}^B(d_P^{(b)}-\bar d_P)^2}.
\end{equation}
Одна перевыборка используется для обоих разбиений в разности, чтобы сохранить зависимость их оценок на одних объектах. При $c\ge0$ задаются допустимое множество и выбранное разбиение:
\begin{equation}\label{eq:adaptive}
 \mathcal E_c=\{P\in\mathcal P:d_P\le c\widehat\sigma_P+\varepsilon\},\qquad
 \Psilh=\arg\max_{P\in\mathcal E_c}I(P),
\end{equation}
где $\varepsilon=10^{-12}$ --- численный допуск. Основное значение $c=1$.

Индекс $\mathrm{silh\_max}$ обозначает максимум silhouette внутри допустимого множества $\mathcal E_c$. Из определения следует $A_L(\Psilh)\le A_L(\Pari)$, однако порядок $Q_T(\Psilh)$ и $Q_T(\Pari)$ может быть любым. Этим объясняется возможность положительного эффекта на скрытой проверочной части без противоречия максимальности~\eqref{eq:leader}.

Правило~\eqref{eq:adaptive} не является тестом равенства качества. Бутстреп условен относительно наблюдённой разметки, а отбор лидера использует ту же выборку. Для множества $\mathcal E_c$ не заявляются доверительное покрытие или известная вероятность ошибки~\citep{efron2014}. При отсутствии классов в $L$ их информация не появляется ни в одной перевыборке.

\subsection{Покрытие и оцениваемый эффект}
Истинное покрытие классов определяется ретроспективно:
\begin{equation}
 \operatorname{cov}(L)=
 \frac{|\{y_i:i\in L\}|}{|\{y_i:i=1,\ldots,n\}|}.
\end{equation}
Оно используется для объяснения поведения методов, но не для вычисления~\eqref{eq:adaptive}. Покрытие классов отличается от доли назначенных объектов $q$ в плотностной кластеризации.

После фиксации решений раскрываются проверочные метки. Основной показатель --- потеря относительно лучшего кандидата того же пула:
\begin{equation}
 R_T(\widehat P)=\max_{P\in\mathcal P}Q_T(P)-Q_T(\widehat P)\ge0.
\end{equation}
Отдельно учитывается абсолютное качество: малая потеря в слабом пуле не означает хорошей кластеризации. Эффект адаптивного выбора относительно внешнего лидера равен
\begin{equation}
 G_T=Q_T(\Psilh)-Q_T(\Pari)
     =R_T(\Pari)-R_T(\Psilh).
\end{equation}
Положительное значение означает улучшение. Для диагностики вводится индикатор переключения
\begin{equation}
 Z_L=\mathbf 1\{\Psilh\ne\Pari\}.
\end{equation}
Если $Z_L=0$, то $G_T=0$. Обратное неверно: разные разбиения могут иметь одинаковый проверочный ARI. Поэтому частота переключений оценивается отдельно от долей улучшений и ухудшений. При положительной вероятности переключения
\begin{equation}
 \mathbb E[G_T]=\Pr(Z_L=1)\,\mathbb E[G_T\mid Z_L=1],
\end{equation}
где ожидание относится к зафиксированной схеме раскрытия меток при заданных данных, пуле и проверочной части. Эта запись разделяет частоту вмешательства и его среднее последствие. Для фиксированного набора исследуются среднее $G_T$ по повторениям разметки и доли улучшений, совпадений и ухудшений. Эти повторы не являются независимыми наборами данных.

Базовые сравнения включают максимум внутреннего индекса без меток, внешний максимум, геометрическое разрешение равенств, фиксированную смесь, смесь с весом, замороженным на отдельной development-серии, фиксированный допуск по частичному ARI и выбор по частоте бутстрепных побед без внутреннего индекса. В коде им соответствуют M0--M7. $\Pari$ и $\Psilh$ соответствуют M1 и M6. Настройка по проверочным меткам подтверждающей серии не допускается.

\subsection{Экспериментальные условия и проверяемые гипотезы}
Основные факторы --- бюджет $m$, схема раскрытия меток и структура данных. Случайный сценарий выбирает объекты равномерно без возвращения. В сценарии отсутствующих классов метки раскрываются только из случайно выбранной половины категорий, хотя все объекты остаются в $X$. В смещённом сценарии половина классов получает больший вес отбора. Бюджеты внутри повторения вложены. Все методы получают одно и то же $L$.

Первый этап исследования --- диагностический анализ сохранённых пилотов. Для каждого повторения сопоставляются внешний лидер и окончательный выбор, их частичный и проверочный ARI, внутренние оценки, размер допустимого множества, отставание и бутстрепный разброс. Сравнение с разрешением только равенств, фиксированным допуском и бутстрепным выбором без внутреннего индекса позволяет отделить эффект геометрии от эффекта самого допуска. Повторные вычисления при разных начальных состояниях генератора бутстрепа оценивают чувствительность решений к конечному $B$.

После диагностики предусматривается расширенная серия с независимыми синтетическими реализациями и несколькими реальными числовыми наборами. При одинаковом бюджете отдельно варьируются покрытие классов и распределение меток между ними. Состав пула и соответствие геометрии целевым классам рассматриваются как самостоятельные факторы. Параметры правила определяются на отдельной настроечной серии и фиксируются до анализа проверочной серии. Существующие пилоты, использованные для разработки правил, не считаются независимым подтверждением их эффективности.

Проверяются следующие предположения: (H1) польза внутреннего индекса зависит от соответствия геометрии целевым классам. (H2) влияние бюджета меток нельзя объяснить без учёта их состава. (H3) при неполном покрытии малый бутстрепный разброс может сопровождаться существенной ошибкой выбора. Основные результаты сравнения должны включать среднюю потерю, частоту и последствия переключений, вероятность ухудшения и величину отрицательных эффектов. Повторы разметки характеризуют вариативность внутри задачи. Переносимость выводов оценивается по независимым задачам. Ни одна из гипотез не предполагает универсального преимущества адаптивного правила.
